The examination is a multiple-choice questionnaire.
A question shows a snippet of code and asks what it computes,
or shows several and asks which is the right one for a stated purpose.
It is marked out of 100, and the pass mark is 60.
Sample questions are published during the course,
so that the format can be practised beforehand.

Two assignments are offered, one for each chapter, and each is marked out of 25.
They are optional, and a student may submit one, both or neither.
Their marks are not added to the questionnaire;
they lower its pass mark instead.
A student who scores $a_{1}$ on the first and $a_{2}$ on the second
passes with a questionnaire score of at least
\begin{equation}\label{eq.passMark}
 60 - a_{1} - a_{2} ,
\end{equation}
so the pass mark falls to 10 out of 100 for a student who scores full marks on both.
